\documentclass[11pt]{article}
\usepackage[utf8]{inputenc}
\usepackage[spanish,es-tabla]{babel}
\usepackage{amsmath,cancel}
\usepackage{amsfonts}
\usepackage{amssymb}
\usepackage{listings}
\usepackage{xcolor}
\usepackage{tabularx}
\usepackage{adjustbox}
\spanishdecimal{.}
\usepackage{graphicx}
\usepackage{float}
\usepackage[lmargin=2.5cm,rmargin=2.5cm,top=2.5cm,
bottom=2cm]{geometry}
\newcommand{\nombre}{Oscar}
\newcommand{\fecha}{3 de mayo del 2026}
\newcommand{\tituloPractica}{}
%%%%%%%%%%%%%%%%%%%%%%%%%%%%%%%%%%%%%%%%%%%%%%
\begin{document}
\begin{titlepage}
\begin{center}
\vspace*{2\baselineskip}
\hrule height 3pt
\vspace*{0.5\baselineskip}
{\Large \textbf{UNIVERSIDAD DE GUADALAJARA}}\\[0.1cm]
{\large \textbf{CENTRO UNIVERSITARIO DE CIENCIAS EXACTAS E INGENIER\'IAS}}
\vspace*{0.5\baselineskip}
\hrule
\vspace*{3\baselineskip}
\includegraphics[width=0.2\textwidth]{udg}
\vspace*{2\baselineskip}\\
{\large\textbf{ Algor\'itmos Bio-Inspirados \\ Msc. Rob\'otica e Inteligencia Artificial}
\vspace*{\baselineskip}\\
{\large Algoritmo 6 --- Particle Swarm Optimization Local Best N-Dimensional \vspace{1cm} \\
\begin{tabular}{p{4cm}p{8cm}}
\hline \hline
Nombre del alumno:  & Oscar Alberto Gallo Garc\'ia \\
C\'odigo del alumno:  & 218357704 \\
Profesor: & Dr. Carlos Alberto Lopez Franco \\
 \hline \hline
\end{tabular}
\vfill
\fecha
\end{center}
\newpage
\end{titlepage}

\tableofcontents
\newpage

\section{Resumen}

Se implement\'o la variante \textit{local best} (lbest) del algoritmo de Optimizaci\'on por Enjambre de Part\'iculas en $n$ dimensiones. A diferencia de la variante \textit{global best} (gbest), en lbest cada part\'icula se orienta \'unicamente por el mejor valor encontrado dentro de su \textbf{vecindario local}, una subconjunto reducido del enjambre definido mediante una topolog\'ia de anillo. Esta restricci\'on en el flujo de informaci\'on preserva mayor diversidad y reduce la convergencia prematura en funciones multimodales. Se evaluaron doce casos: minimizaci\'on y maximizaci\'on en 2~dimensiones con gr\'aficas, y minimizaci\'on y maximizaci\'on en 10~dimensiones.

\section{Metodolog\'ia}

\subsection{Diferencia con gbest}

En la variante gbest, todas las part\'iculas comparten y se atraen hacia un \'unico $\hat{\mathbf{y}}$ global. En la variante lbest, cada part\'icula $i$ tiene su propio \textit{local best} $\hat{\mathbf{y}}_i$, definido como la mejor posici\'on personal encontrada dentro de su vecindario $\mathcal{N}_i$.

\begin{table}[H]
\centering
\begin{tabular}{lll}
\hline\hline
Caracter\'istica & gbest & lbest \\
\hline
Vecindario     & enjambre completo        & subconjunto de $k$ part\'iculas \\
$\hat{\mathbf{y}}$ & \'unico para todos   & distinto para cada part\'icula \\
Velocidad de convergencia & r\'apida         & m\'as lenta \\
Diversidad     & baja                     & alta \\
Riesgo de \'optimo local & mayor            & menor \\
\hline\hline
\end{tabular}
\caption{Comparaci\'on entre las variantes gbest y lbest.}
\end{table}

\subsection{Topolog\'ia de anillo}

Se utiliza una topolog\'ia de \textbf{anillo} con tama\~no de vecindario $k = 3$: cada part\'icula $i$ tiene como vecinos a las part\'iculas $i-1$, $i$ y $i+1$ (con \textit{wrap-around} en los extremos):
\[
\mathcal{N}_i = \{(i-1)\bmod n_s,\; i,\; (i+1)\bmod n_s\}
\]

El \textit{local best} para la part\'icula $i$ es:
\[
\hat{\mathbf{y}}_i(t) \in \{\mathbf{y}_j(t) : j \in \mathcal{N}_i\}
\;\Big|\;
f\!\left(\hat{\mathbf{y}}_i(t)\right) = \min_{j \in \mathcal{N}_i} f\!\left(\mathbf{y}_j(t)\right)
\]

\subsection{Actualizaci\'on de velocidad}

La ecuaci\'on de velocidad conserva la misma estructura que gbest, sustituyendo $\hat{\mathbf{y}}$ global por $\hat{\mathbf{y}}_i$ local:
\[
v_{ij}(t+1) = w\,v_{ij}(t)
  + c_1\,r_{1j}(t)\bigl[y_{ij}(t) - x_{ij}(t)\bigr]
  + c_2\,r_{2j}(t)\bigl[\hat{y}_{ij}(t) - x_{ij}(t)\bigr]
\]

La posici\'on se actualiza igual que en gbest:
\[
\mathbf{x}_i(t+1) = \mathbf{x}_i(t) + \mathbf{v}_i(t+1)
\]

\subsection{Algoritmo lbest}

\begin{enumerate}
    \item Inicializar el enjambre: $\mathbf{x}_i \sim \mathcal{U}(\ell_{\min},\ell_{\max})$,\quad $\mathbf{v}_i \sim \mathcal{U}(-\Delta,\Delta)$,\quad $\mathbf{y}_i = \mathbf{x}_i$.
    \item Evaluar $f(\mathbf{y}_i)$ para toda $i$.
    \item Repetir hasta la condici\'on de paro:
    \begin{enumerate}
        \item Para cada part\'icula $i$: si $f(\mathbf{x}_i) < f(\mathbf{y}_i)$, actualizar $\mathbf{y}_i \leftarrow \mathbf{x}_i$.
        \item Para cada part\'icula $i$:
        \begin{itemize}
            \item Determinar $\hat{\mathbf{y}}_i \leftarrow \arg\min_{j \in \mathcal{N}_i} f(\mathbf{y}_j)$.
            \item Actualizar $\mathbf{v}_i$ y $\mathbf{x}_i$; proyectar $\mathbf{x}_i$ al dominio $[\ell_{\min},\ell_{\max}]$.
        \end{itemize}
    \end{enumerate}
    \item Retornar $\arg\min_{i}\, f(\mathbf{y}_i)$.
\end{enumerate}

\subsection{Par\'ametros}

\begin{table}[H]
\centering
\begin{tabular}{lcc}
\hline\hline
Par\'ametro & 2D & 10D \\
\hline
$w$                       & 0.7  & 0.7  \\
$c_1$                     & 1.5  & 1.5  \\
$c_2$                     & 1.5  & 1.5  \\
$k$ (tama\~no vecindario) & 3    & 3    \\
$n_s$ (tama\~no enjambre) & 50   & 100  \\
Iteraciones m\'aximas     & 200  & 1000 \\
\hline\hline
\end{tabular}
\caption{Par\'ametros utilizados en todos los experimentos.}
\end{table}

\section{Resultados}

\subsection{Minimizaci\'on 2D}

\subsubsection{Funci\'on Esfera}

M\'inimo global en $\mathbf{x}^* = (0,0)$. Mejor soluci\'on: $(-0,\;0)$,\quad $f = 2.71\times10^{-21}$.

\begin{figure}[H]
\centering
\includegraphics[width=14cm]{figs/PSO_local_min_esfera.png}
\caption{Minimizaci\'on de la funci\'on Esfera con lbest PSO. La estrella roja indica la mejor soluci\'on del enjambre.}
\end{figure}

\subsubsection{Funci\'on de Rastrigin}

M\'inimo global en $\mathbf{x}^* = (0,0)$. Mejor soluci\'on: $(-0,\;-0)$,\quad $f = 0.0$.

\begin{figure}[H]
\centering
\includegraphics[width=14cm]{figs/PSO_local_min_rastrigin.png}
\caption{Minimizaci\'on de Rastrigin. La diversidad mantenida por la topolog\'ia de anillo ayuda a escapar de los m\'inimos locales peri\'odicos.}
\end{figure}

\subsubsection{Funci\'on de Rosenbrock}

M\'inimo global en $\mathbf{x}^* = (1,1)$. Mejor soluci\'on: $(0.99962,\;0.99923)$,\quad $f = 1.49\times10^{-7}$.

\begin{figure}[H]
\centering
\includegraphics[width=14cm]{figs/PSO_local_min_rosenbrock.png}
\caption{Minimizaci\'on de Rosenbrock. El enjambre navega el valle curvo y converge cerca del m\'inimo global.}
\end{figure}

\subsection{Maximizaci\'on 2D}

\subsubsection{Esfera Negativa}

M\'aximo global en $(0,0)$. Mejor soluci\'on: $(0,\;0)$,\quad $f = -8.50\times10^{-23}$.

\begin{figure}[H]
\centering
\includegraphics[width=14cm]{figs/PSO_local_max_neg_esfera.png}
\caption{Maximizaci\'on de la esfera negativa. El enjambre converge al v\'ertice de la par\'abola invertida.}
\end{figure}

\subsubsection{Funci\'on Seno-Coseno}

M\'aximo global de $1.0$. Mejor soluci\'on: $(-1.5708,\;3.1416)$,\quad $f = 1.0$.

\begin{figure}[H]
\centering
\includegraphics[width=14cm]{figs/PSO_local_max_seno_coseno.png}
\caption{Maximizaci\'on de $\sin(x_1)\cos(x_2)$. El lbest localiza uno de los m\'aximos id\'enticos con valor exacto $f=1$.}
\end{figure}

\subsubsection{Funci\'on Easom}

M\'aximo global en $(\pi,\pi)$ con $f = 1.0$. Mejor soluci\'on: $(3.1416,\;3.1416)$,\quad $f = 1.0$.

\begin{figure}[H]
\centering
\includegraphics[width=14cm]{figs/PSO_local_max_easom.png}
\caption{Maximizaci\'on de Easom. La topolog\'ia de anillo preserva diversidad en el dominio casi plano hasta localizar el pico estrecho en $(\pi,\pi)$.}
\end{figure}

\subsection{Minimizaci\'on N-Dimensional (10D)}

\begin{table}[H]
\centering
\begin{tabular}{lp{8.5cm}r}
\hline\hline
Funci\'on & Mejor soluci\'on & $f(\hat{\mathbf{y}})$ \\
\hline
Esfera    & $[0,\;{-0},\;{-0},\;{-0},\;0,\;0,\;0,\;0,\;{-0},\;0]$                              & $4.07\times10^{-32}$ \\[3pt]
Rastrigin & $[0,\;0,\;{-0},\;0.995,\;{-0},\;0.995,\;0,\;1.990,\;0,\;{-0}]$                    & $5.97$ \\[3pt]
Ackley    & $[{-0},\;{-0},\;0,\;0,\;0,\;0,\;0,\;0,\;0,\;0]$                                   & $2.89\times10^{-14}$ \\
\hline\hline
\end{tabular}
\caption{Resultados de minimizaci\'on en 10 dimensiones.}
\end{table}

\subsection{Maximizaci\'on N-Dimensional (10D)}

\begin{table}[H]
\centering
\begin{tabular}{lp{8.5cm}r}
\hline\hline
Funci\'on & Mejor soluci\'on & $f(\hat{\mathbf{y}})$ \\
\hline
$-$Esfera    & $[0,\;0,\;0,\;{-0},\;0,\;{-0},\;0,\;0,\;0,\;0]$                               & $-1.60\times10^{-32}$ \\[3pt]
$-$Rastrigin & $[0.995,\;0,\;0,\;{-0},\;{-0},\;{-0},\;{-0.995},\;{-0},\;{-0},\;{-0}]$       & $-1.99$ \\[3pt]
$-$Ackley    & $[{-0},\;0,\;{-0},\;{-0},\;0,\;0,\;{-0},\;{-0},\;{-0},\;0]$                  & $-1.11\times10^{-14}$ \\
\hline\hline
\end{tabular}
\caption{Resultados de maximizaci\'on en 10 dimensiones.}
\end{table}

\subsection{Comparaci\'on lbest vs.\ gbest en 10D}

\begin{table}[H]
\centering
\begin{tabular}{llrr}
\hline\hline
Modo & Funci\'on & gbest & lbest \\
\hline
Minimizaci\'on & Esfera    & $6.31\times10^{-69}$ & $4.07\times10^{-32}$ \\
               & Rastrigin & $6.96$               & $\mathbf{5.97}$ \\
               & Ackley    & $4.00\times10^{-15}$ & $2.89\times10^{-14}$ \\
\hline
Maximizaci\'on & $-$Esfera    & $-1.08\times10^{-69}$ & $-1.60\times10^{-32}$ \\
               & $-$Rastrigin & $-5.97$               & $\mathbf{-1.99}$ \\
               & $-$Ackley    & $-4.00\times10^{-15}$ & $-1.11\times10^{-14}$ \\
\hline\hline
\end{tabular}
\caption{Comparaci\'on de $f(\hat{\mathbf{y}})$ entre las variantes gbest y lbest (mismos par\'ametros y semilla). En negrita el mejor resultado en Rastrigin, funci\'on multimodal.}
\end{table}

\section{Conclusi\'on}

El lbest PSO encontr\'o el \'optimo global en las funciones unimodales (Esfera, Ackley, Rosenbrock) con precisi\'on comparable a gbest, y obtuvo un mejor resultado en Rastrigin 10D ($f=5.97$ frente a $f=6.96$ de gbest), lo que refleja la ventaja de la topolog\'ia de anillo en funciones altamente multimodales.\\\\
La raz\'on es que en gbest la informaci\'on del mejor global se propaga instant\'aneamente a todo el enjambre, favoreciendo una convergencia r\'apida pero prematura. En lbest la informaci\'on se difunde lentamente a trav\'es de los vecinos, lo que preserva la diversidad y permite a regiones del enjambre explorar zonas distintas del espacio antes de comprometerse con un \'optimo.\\\\
La desventaja es que lbest converge m\'as lentamente y, con los mismos par\'ametros, la precisi\'on en funciones unimodales es ligeramente menor que gbest (p.~ej.\ Esfera 10D: $4\times10^{-32}$ vs $6\times10^{-69}$). La elecci\'on entre variantes depende del problema: gbest es preferible en espacios de b\'usqueda unimodales; lbest es m\'as robusto frente a m\'ultiples \'optimos locales.

\end{document}
